\renewcommand{\Anno}{2023}

\renewcommand{\NumeroProblema}{2}

\renewcommand{\Testo}{
La Scuola Sant’Anna intende acquistare in data 1/1/2024 due hard disk dello stesso tipo. Ciascuno di questi hard disk funziona bene per tutto un anno con probabilità $p$, e di conseguenza si guasta nel corso di un anno con probabilità $1-p$.

\begin{enumerate}
\renewcommand{\labelenumi}{(\alph{enumi})}

\item  Determinare la probabilità che \emph{entrambi} gli hard disk si guastino nel corso dei primi tre anni successivi all'acquisto (cioè entro il 31/12/2026).

\item  Determinare la probabilità che entrambi gli hard disk si guastino \emph{nello stesso anno} (ma \emph{non} necessariamente entro il 31/12/2026).

\item  Supponiamo ora che il secondo hard disk sia di qualità inferiore rispetto al primo, per cui la sua probabilità di funzionare bene per tutto un anno è $p/2$ invece di $p$.

Determinare la probabilità che il secondo hard disk  si guasti in un anno strettamente maggiore rispetto all'anno in cui si guasta il primo (quello di qualità superiore).

\end{enumerate}
}

\renewcommand{\Risposta}{
% SSSUP-BOZZA-AI: Risposta; autore=ChatGPT; data=2026-11-08; stato=DA-VERIFICARE
% SSSUP-PROVENIENZA: AI: elaborazione nuova dal testo
Assumendo guasti indipendenti e probabilità di sopravvivenza annuale costante, condizionata al funzionamento iniziale:
\begin{enumerate}
\item[(a)] $(1-p^3)^2$.
\item[(b)] $(1-p)/(1+p)$ per $0\leq p<1$; per $p=1$ la probabilità è zero.
\item[(c)] $p(1-p)/(2-p^2)$.
\end{enumerate}
Senza tali ipotesi probabilistiche, le sole informazioni marginali del testo non determinano univocamente le risposte.
}

\renewcommand{\Aiutino}{
% SSSUP-BOZZA-AI: Aiutino; autore=ChatGPT; data=2026-11-08; stato=DA-VERIFICARE
% SSSUP-PROVENIENZA: AI: elaborazione nuova dal testo
Modella l'anno del primo guasto con una variabile geometrica: $\Pr(T=k)=p^{k-1}(1-p)$. Nei punti (b) e (c) somma sugli anni possibili e usa una serie geometrica. Distingui il caso $p=1$, in cui il primo disco non si guasta mai.
}

\renewcommand{\SoluzioneUfficiale}{}

\renewcommand{\SoluzioniAlternative}{}

\renewcommand{\Osservazioni}{}

\renewcommand{\Commenti}{}

% Aggiunta redazionale: la soluzione preesistente non viene modificata.
\providecommand{\SoluzioneRedazionale}{}
\renewcommand{\SoluzioneRedazionale}{
% SSSUP-BOZZA-AI: SoluzioneRedazionale; autore=ChatGPT; data=2026-11-08; stato=DA-VERIFICARE
% SSSUP-PROVENIENZA: AI: elaborazione nuova dal testo
Esplicitiamo il modello probabilistico usuale: i dischi sono indipendenti e, per ciascun disco ancora funzionante, la probabilità di sopravvivere al successivo anno è sempre la stessa. Con parametro $p<1$, l'anno $T$ del primo guasto ha quindi distribuzione
\[
\Pr(T=k)=p^{k-1}(1-p),\qquad k=1,2,\ldots.
\]
Queste ipotesi sono necessarie: il testo non specifica completamente le dipendenze fra i guasti.

\paragraph{Punto (a).}
Un disco sopravvive ai primi tre anni con probabilità $p^3$, quindi si guasta entro quel termine con probabilità $1-p^3$. Per indipendenza, la probabilità che ciò avvenga per entrambi è $(1-p^3)^2$.

\paragraph{Punto (b).}
Per $p<1$,
\[
\Pr(T_1=T_2)=\sum_{k=1}^{\infty}(1-p)^2p^{2(k-1)}
=\frac{(1-p)^2}{1-p^2}=\frac{1-p}{1+p}.
\]
Per $p=1$ non avviene alcun guasto e la probabilità richiesta è zero; non si considera ``guastarsi nello stesso anno'' il non guastarsi mai.

\paragraph{Punto (c).}
Se il primo disco si guasta nell'anno $k$, il secondo deve sopravvivere a tutti i primi $k$ anni, evento di probabilità $(p/2)^k$. Pertanto
\[
\Pr(T_2>T_1)=\sum_{k=1}^{\infty}(1-p)p^{k-1}\left(\frac p2\right)^k
=\frac{p(1-p)}{2-p^2}.
\]
La formula vale anche agli estremi: per $p=0$ entrambi si guastano nel primo anno; per $p=1$ il primo non si guasta mai e il secondo non può guastarsi in un anno successivo al suo.
}
